% Hafta 2 — Karşı önlemler: dört katman, dört soru
% Derleme: py -3.12 tools/latex/derle.py docs/week-2/assets/h02-10-karsi-onlem-katmanlari.tex
\documentclass[tikz,border=8pt]{standalone}
\usepackage{cen429-cizim}
\begin{document}
\begin{tikzpicture}[node distance=6mm and 6mm]
  \node[baslik] (b) at (0,0) {Karşı önlemler: dört katman, dört soru};
  \node[altbaslik] at (b.south west) {her katman ötekinin kaçırdığını yakalar};

  \node[kutu=48mm, minimum height=13mm, below=13mm of b.south west, anchor=north west, xshift=6mm] (k1)
    {\kb{İmza tabanlı}};
  \node[font=\sffamily\small, text=soluk, text width=110mm, align=left, right=6mm of k1] (a1)
    {``bu dosyayı daha önce gördüm mü?'' --- bilinen zararlıyı yakalar};

  \node[kutu=48mm, minimum height=13mm, below=5mm of k1] (k2) {\kb{Sezgisel (heuristic)}};
  \node[font=\sffamily\small, text=soluk, text width=110mm, align=left, right=6mm of k2] (a2)
    {``şüpheli görünüyor mu?'' --- benzerini yakalar};

  \node[iyikutu=48mm, minimum height=13mm, below=5mm of k2] (k3) {\kb{Davranış tabanlı}};
  \node[font=\sffamily\small, text=soluk, text width=110mm, align=left, right=6mm of k3] (a3)
    {``ne YAPIYOR?'' --- yeni zararlıyı çalışırken yakalar};

  \node[iyikutu=48mm, minimum height=13mm, below=5mm of k3] (k4) {\kb{Bütünlük izleme}};
  \node[font=\sffamily\small, text=soluk, text width=110mm, align=left, right=6mm of k4] (a4)
    {``değişti mi?'' --- beyaz liste, kurcalamayı yakalar};

  \node[dip, text width=185mm, below=10mm of k4.south, anchor=north]
    {Tek katman yetmez: imza yeniyi, davranış sessizi kaçırır.};
\end{tikzpicture}
\end{document}
